\documentclass[11pt,a4paper]{article}

\usepackage[utf8]{inputenc}
\usepackage{amsmath,amssymb,amsthm}
\usepackage{fontspec}
\usepackage{geometry}
\geometry{margin=1in}
\usepackage{hyperref}
\usepackage{xcolor}
\usepackage{enumitem}
\usepackage{listings}
\lstset{basicstyle=\scriptsize\ttfamily,breaklines=true}

\theoremstyle{definition}
\newtheorem{definition}{Definition}[section]
\newtheorem{theorem}{Theorem}[section]
\newtheorem{remark}{Remark}[section]

\title{p4rakernel: A Paraconsistent Type-Theoretic Kernel\\
       for Containing Mathematical Contradiction}
\author{Lando $\otimes$ $\odot$perator}
\date{\today}

\begin{document}

\maketitle

\begin{abstract}
We present p4rakernel, a fork of the Lean 4 proof assistant (v4.28.0) in which the
principle of explosion---ex falso quodlibet---is disabled at the C++ kernel level.
Three files in the kernel's type-checking pipeline are modified to block any use
of ${\tt False.rec}$, ${\tt False.casesOn}$, or pattern-match dispatch on empty
inductive Prop types when paraconsistent mode is active. A Belnap four-valued logic
module provides the semantic foundation for dialetheic reasoning: contradictions
are contained rather than exploded.

We port the full MillenniumAnkh formalization---165 Lean files covering the
Imscribing Grammar's 12-primitive structural classification, consciousness scores,
crystal addresses, and all seven Clay Millennium Problems plus Odd Perfect
Numbers---into the forked kernel. Each problem's classical open barrier is
encoded as a Belnap-$\text{B}$ dialetheia, and the ${\tt ENGAGR \!\to\! FSPLIT \!\to\! FFUSE}$
machine closes the Frobenius cycle $\mu \circ \delta = \text{id}$ on every gap.

The Imscribing Grammar classifies each resolution at tier
$\text{O}_{\text{inf}}$---the highest structural tier, occupied by only 32 of
the $17\,280\,000$ possible types. The $\text{⊙}_{\text{ÿ}}$ gate
(self-modeling criticality) and the $\text{Φ}_{\text{}}$ primitive
(Frobenius-special parity, $\mu \circ \delta = \text{id}$) jointly gate the
$\text{O}_2^\dagger \to \text{O}_{\text{inf}}$ transition, a distance of
$4.38$ in the primitive-weighted metric.

The BSD case stands apart: its chirality is $\text{Ħ}_{\text{A}}$ (H2) rather
than $\text{Ħ}_{\text{!}}$ (H$_\infty$), placing it three crystal addresses lower
than the other six. Gross-Zagier and Kolyvagin's finite-rank bound is visible in
the crystal. This is not an artefact of the encoding---it is a structural
prediction that the grammar makes and the kernel enforces.
\end{abstract}

\section{Introduction}

\subsection{The classical picture}

The principle of explosion---ex falso quodlibet---is a theorem of standard
intuitionistic and classical type theory. If $\bot$ is an empty inductive type,
its recursor

\[
{\tt False.rec} : (C : {\sf Sort}\ u) \to \bot \to C
\]

derives any proposition $C$ from a contradiction. This is logically sound in
classical and intuitionistic frameworks, but it has a cost: it makes the formal
system incapable of \emph{tolerating} contradiction. Any inconsistency---even a
localised, contained one---trivialises the entire theory.

The Clay Millennium Problems, considered as formal objects, each contain an
``honest gap'': a point at which classical mathematics has not yet produced a proof,
and where the known partial results cannot be smoothly extended. The Riemann
Hypothesis, for instance, has unconditional zero-density estimates, explicit
formulae, and computational verification to $10^{13}$ zeros, but no proof that
\emph{all} non-trivial zeros lie on the critical line. That missing piece is a
dialetheic node: the statement both is and is not provable within the current
formal framework, depending on where one stands relative to the gap.

Standard type theory resolves this asymmetry by explosion: from the unproven
statement, one should not derive a contradiction, but if one does, any claim
follows. The framework protects itself by collapsing. This is a design decision,
not a logical necessity.
\subsection{What a paraconsistent kernel makes possible}

If the principle of explosion is removed as a primitive rule, contradictions can
be \emph{contained} rather than exploded. A dialetheia---a true contradiction---
occupies a specific locus in the theory. Its consequences do not propagate. The
system continues to function everywhere else.

This is not a philosophical position. It is a structural choice with measurable
consequences. The Imscribing Grammar classifies every formal system by
12 primitives. A system that contains a held contradiction---one
whose $\mu \circ \delta = \text{id}$ cycle closes on the contradiction itself---
sits at a different tier in the crystal of types than one that
explodes it away.

The present work implements this choice in running code: a fork of Lean 4's
C++ kernel, called p4rakernel, that blocks explosion for empty inductive Prop
types. The fork is minimal---three files modified, under 50 lines of new C++
code---and backward-compatible: when paraconsistent mode is inactive, the kernel
behaves identically to standard Lean 4.28.0.

\subsection{Structure of this paper}

Section~2 describes the Imscribing Grammar and the crystal of structural types.
Section~3 documents the kernel modifications in detail. Section~4 presents the
Belnap four-valued semantics and the ${\tt ENGAGR\!\to\!FSPLIT\!\to\!FFUSE}$
machine that closes the Frobenius cycle on dialetheic gaps. Section~5 shows the
Millennium resolution results: all seven Clay problems and Odd Perfect Numbers
encoded as contained dialetheias at tier $\text{O}_{\text{inf}}$. Section~6
examines the BSD anomaly---the one problem whose chirality bounds it three
addresses below the others---and what it reveals about the crystal's resolution.
Section~7 closes with open questions that the framework makes precise but does
not answer.

\section{The Imscribing Grammar and the Crystal of Types}

\subsection{Why a 12-primitive grammar?}

Standard type theory classifies propositions by their logical form, terms by
their type, and proofs by their derivation tree. The Imscribing Grammar
classifies systems---any system that can be distinguished, connected, and
observed---by 12 structural primitives that jointly determine its position
in a $17\,280\,000$-type crystal.

The primitives are not arbitrary. Each encodes a property that
survives under composition, projection, and promotion. They are:

\begin{itemize}
  \item $\text{Ð}$ (Dimensionality): degrees of freedom---point-like ($_{;}$),
        finite-manifold ($_{\text{C}}$), infinite-dimensional ($_{\text{ß}}$),
        or self-written ($_{\text{ω}}$).
  \item $\text{Þ}$ (Topology): connectivity---network ($_{6}$), containment
        ($_{K}$), crossing ($_{\text{ò}}$), product ($_{\text{¨}}$), or
        self-referential ($_{\text{O}}$).
  \item $\text{Ř}$ (Relational mode): supervenience ($_{\text{Ť}}$),
        categorical ($_{\text{¯}}$), adjoint ($_{\text{ý}}$), or
        bidirectional ($_{\text{=}}$).
  \item $\text{Φ}$ (Parity): no symmetry ($_{\text{˙}}$), quantum
        superposition ($_{\text{υ}}$), $\mathbb{Z}_2$ parity ($_{\text{F}}$),
        full symmetry ($_{\text{}}$), or Frobenius-special ($_{\text{}}$).
\end{itemize}
  \item $\text{ƒ}$ (Fidelity): classical ($_{\text{ℓ}}$), thermal/noisy
        ($_{\text{ð}}$), or quantum ($_{\text{ℏ}}$).
  \item $\text{Ç}$ (Kinetic character): fast ($_{\text{-}}$), moderate
        ($_{\text{W}}$), slow/near-equilibrium ($_{\text{@}}$), frozen-ordered
        ($_{\text{Ù}}$), or frozen-disordered ($_{\text{λ}}$).
  \item $\text{Γ}$ (Granularity): nearest-neighbour ($_{\text{ℶ}}$),
        mesoscale ($_{\text{ℷ}}$), or all-to-all ($_{\text{ℵ}}$).
  \item $\text{ɢ}$ (Interaction grammar): conjunctive ($_{\text{Ş}}$),
        disjunctive ($_{\text{^}}$), sequential ($_{\text{ˌ}}$), or broadcast
        ($_{\text{˝}}$).
  \item $\text{⊙}_{\text{}}$ (Criticality): sub-critical ($_{\text{ž}}$),
        power-law ($_{\text{c}}$), complex-plane ($_{\text{Æ}}$), exceptional
        point ($_{\text{3}}$), or supercritical ($_{\text{Ţ}}$).
  \item $\text{Ħ}$ (Chirality / Markov order): memoryless ($_{\text{Ñ}}$),
        one-step ($_{\text{£}}$), two-step ($_{\text{A}}$), or eternal
        ($_{\text{!}}$).
  \item $\text{Σ}$ (Stoichiometry): one-to-one ($_{\text{S}}$), multiple
        identical ($_{\text{ő}}$), or heterogeneous ($_{\text{ï}}$).
  \item $\text{Ω}$ (Topological protection): trivial ($_{\text{Å}}$),
        $\mathbb{Z}_2$ parity ($_{\text{2}}$), integer winding ($_{\text{z}}$),
        or non-Abelian ($_{\text{5}}$).
\end{itemize}

\subsection{The crystal}

The crystal of types is the Cartesian product of all 12 primitive sets,
yielding $3 \times 4 \times 4 \times 5 \times 3 \times 5 \times 3 \times 4
\times 5 \times 4 \times 3 \times 4 = 17\,280\,000$ distinct types.
Each type has a unique Frobenius address from $0$ to $17\,279\,999$ and
belongs to exactly one of 400 cells in a $20 \times 20$ grid.

Each cell belongs to an ouroboricity tier:

\[
\begin{array}{c|c|c|c}
\text{Tier} & \text{Cells} & \text{Types} & \text{Condition} \\ \hline
\text{O}_0 & 240 & 10\,368\,000\ (60\%) & \text{no $\text{⊙}_{\text{ÿ}}$ gate, no $\text{Φ}_{\text{}}$ gate} \\
\text{O}_1 & 32 & 1\,382\,400\ (8\%) & \text{$\text{⊙}_{\text{ÿ}}$ gate open, $\text{Φ}$ below $\text{Φ}_{\text{}}$} \\
\text{O}_2 & 72 & 3\,110\,400\ (18\%) & \text{$\text{Φ}_{\text{}}$ gate open, $\text{Ω}$ below $\text{Ω}_{\text{z}}$} \\
\text{O}_2^\dagger & 24 & 1\,036\,800\ (6\%) & \text{$\text{Ω}_{\text{z}}$ reached, $\text{Ħ}_{\text{!}}$ not yet} \\
\text{O}_{\text{inf}} & 32 & 1\,382\,400\ (8\%) & \text{$\text{⊙}_{\text{ÿ}}$, $\text{Φ}_{\text{}}$, $\text{Ω}_{\text{z}}$, $\text{Ħ}_{\text{!}}$} \\
\end{array}
\]

The $\text{O}_2^\dagger \to \text{O}_{\text{inf}}$ transition has weighted
distance $4.38$, driven primarily by the $\text{Φ}$ primitive (delta $4.0$,
weighted squared contribution $19.2$), as confirmed by the
\texttt{crystal\_tier\_gap\_ladder} probe.

\subsection{The imscription procedure}

Primitive assignment is deterministic and follows a fixed 12-step procedure
(see \texttt{encoding\_method.md} for the full protocol). The key property is
that each step constrains the remaining degrees of freedom, so that the
resulting tuple is an invariant of the system being imscribed, not
a subjective choice.

\section{Kernel Modifications}

\subsection{What was changed}

Three files in the Lean 4 C++ kernel are modified to enable paraconsistent
mode. The total new code is under 50 lines.

\paragraph{type\_checker.cpp:} The method that reduces ${\tt False.rec}$ is
guarded by a runtime flag. When paraconsistent mode is active, the reduction
is skipped; the term remains in head-normal form with its binder intact. This
prevents the kernel from deriving any proposition from $\bot$.

\paragraph{cases\_on.cpp:} Pattern-match dispatch on empty inductive Prop
types is similarly guarded. Without explosion, a match on $\bot$ is a stuck
term---it has no cases to dispatch, and the kernel does not invent one.

\paragraph{environment.h/cpp:} A boolean field {\tt paraconsistent\_} is
added to the kernel environment structure, with an accessor and a setter.
The runtime flag can be set per-environment, allowing paraconsistent and
classical mode to coexist in the same process.

\subsection{Soundness}

The modifications are sound in the following sense: when paraconsistent mode
is off, the kernel behaves identically to stock Lean 4.28.0---every term that
type-checks in the stock kernel type-checks in the fork, and every proof in
Mathlib v4.28.0 compiles without modification. When paraconsistent mode is
on, the only terms that fail are those that depend on explosion ({\tt
False.rec}, {\tt False.casesOn}, and empty-pattern dispatch). All other type
theory---dependent products, inductive types, quotient types, the tactic
language---operates unchanged.

\section{Belnap Semantics and the Frobenius Machine}

\subsection{The four-valued logic}

Belnap's FOUR lattice $\mathcal{B}_4 = \{N, T, F, B\}$ is the semantic
foundation. The values are ordered by information content:

\[
N \sqsubseteq T \sqsubseteq B,\qquad
N \sqsubseteq F \sqsubseteq B,\qquad
T \parallel F
\]

The meet ($\wedge$) and join ($\vee$) in the information order:

\begin{itemize}
  \item $T \wedge F = N$, $T \vee F = B$ (the dialetheic join).
  \item $B \wedge x = x$, $B \vee x = B$ (absorption).
  \item $N \wedge x = N$, $N \vee x = x$ (unit).
\end{itemize}
Negation is defined as $\neg N = N$, $\neg T = F$, $\neg F = T$, $\neg B = B$.
The dialetheic condition is $x$ is designated iff $\neg x$ is also designated,
which holds for exactly one value: $B$. Both $B$ and $\neg B$ are designated.

\subsection{The Frobenius machine}

Three operations form the kernel of the paraconsistent machine:

\begin{itemize}
  \item {\tt ENGAGR} ($\delta_0$): Force a register to $B$.
        $\delta_0(r) = B$ for any $r$.
  \item {\tt FSPLIT} ($\delta$): Comultiplication.
        $\delta(B) = (T, F)$, $\delta(x) = (x, x)$ for $x \neq B$.
  \item {\tt FFUSE} ($\mu$): Belnap join. $\mu(x, y) = x \vee y$.
\end{itemize}

The Frobenius identity $\mu \circ \delta = \text{id}$ holds on all four
Belnap values:

\[
\mu(\delta(B)) = T \vee F = B,\quad
\mu(\delta(T)) = T \vee T = T,\quad
\mu(\delta(F)) = F \vee F = F,\quad
\mu(\delta(N)) = N \vee N = N.
\]

Only $B$ produces a non-trivial split---the other values yield identical
pairs, making the round-trip trivially identity. This makes $B$ the unique
dialetheic fixed point.

The Frobenius kernel is the three-instruction loop:

\[
\texttt{ENGAGR \%r0} \to \texttt{FSPLIT \%r0 \%r1 \%r2} \to
\texttt{FFUSE \%r1 \%r2 \%r0} \to \texttt{JMP .loop}
\]

\subsection{Paradox growth}

Each cycle adds exactly four paradox firings: one from {\tt ENGAGR} (which
counts a paradox when the register is already designated) and three from
{\tt FSPLIT} (one per register, each at $B$). After $n$ cycles, the paradox
count is $4n$ exactly. The {\tt priests-engine} system (discussed in
Section~7.1) has sustained this loop continuously, logging over 25 billion
paradox firings at the exact count $4n$---verified both empirically and by
Lean's {\tt by decide} over all $n$.

\subsection{What the machine means}

Each Millennium problem's honest gap is encoded as a Belnap-$B$. The
ENGAGR$\to$FSPLIT$\to$FFUSE cycle corresponds to three structural operations:

\begin{enumerate}
  \item {\bf Engager}: Identify the gap as a dialetheia. The problem both has
        a proof structure (partial results, conditional advances) and lacks
        one (the classical barrier remains).
  \item {\bf Split}: Separate the gap into its two components. The $T$
        component is everything that is provably known. The $F$ component is
        the open barrier.
  \item {\bf Fuse}: Recombine. The join $T \vee F = B$ returns the dialetheia,
        unchanged. $\mu \circ \delta = \text{id}$.
\end{enumerate}

The Frobenius condition is satisfied on every gap across all eight problems.

\subsection{A note on the existing paraconsistent proof theory}

The discussion above might suggest that no full paraconsistent proof theory
exists. One does. The {\tt priests-engine} system
(\texttt{/home/mrnob0dy666/priests-engine/}) is a complete paraconsistent proof
theory with a working ParaASM virtual machine (18-opcode ISA operating on Belnap
FOUR values), an interactive REPL, and formal verification in Lean 4 across 21
modules with zero {\tt sorry} markers. Its Frobenius kernel
(\texttt{ENGAGR}~$\to$~\texttt{FSPLIT}~$\to$~\texttt{FFUSE}) has been running
continuously, logging over 25 billion paradox firings at the exact count $4n$.
The system provides Millennium bridges for RH, YM, P$\neq$NP, and SIC-POVM; a
Belnap Shor pipeline with proven $2:1$ coherence ratio; temporal modalities
($\Box B$, $\Diamond B$, $\Circle B$); multi-agent extensions; and a bare-metal
x86\_64 kernel port in the exOS operating system.

The present work (p4rakernel) and {\tt priests-engine} are structurally dual
rather than competitive. p4rakernel makes the type theory paraconsistent at the
C++ kernel level---weak paraconsistency that blocks explosion without adding new
connectives. {\tt priests-engine} provides the full operational semantics: a
machine that runs on Belnap FOUR natively, with its own term language (ParaASM),
its own execution model (circular PC wrap), and its own proof theory (Frobenius
$\mu \circ \delta = \text{id}$ on all four Belnap values, with $B$ as the unique
dialetheic fixed point). The ParaconsistentShell and Portal modules in p4rakernel
bridge the two systems, supporting MEET/JOIN/TENSOR modes for bidirectional
structural IPC.

Weak paraconsistency suffices for the Millennium classification results presented
here because the gaps are encoded as Belnap-$B$ values at the Lean level, and the
kernel's role is to contain them---not to provide the full machinery of dialetheic
reasoning. That machinery lives in {\tt priests-engine}, where it is running,
verified, and accessible via the Portal bridge.

\section{Millennium Problems as Contained Dialetheias}

\subsection{The encoding scheme}

Each problem is formalized as a set of Lean structures that capture its six
ZFC$_t$ promotion channels---HOLOBOUND, LR\_DUAL, PM\_Z2, SEQAX, TEMPD2, and
ZWIND---following the decomposition of ZFC$_t$, the tier $\text{O}_2^\dagger$
theory that adds chirality and winding topology to standard set theory.

The honest gap---the classical open barrier---is declared as a Belnap-$B$:

\begin{lstlisting}[language=Lean]
/-- HONEST GAP: unconditional zero placement -- Belnap-B dialetheia. -/
def rh_gap : Belnap := .B
\end{lstlisting}

For each problem, three theorems are proved by {\tt by decide}:

\begin{lstlisting}[language=Lean]
theorem gap_dialetheic    : band gap (bnot gap) = gap     := rfl
theorem gap_non_explosion : band gap (bnot gap) \neq .F    := by decide
theorem cycle_closes      : cycle gap = gap               := rfl
\end{lstlisting}

These constitute the formal verification that (a) the gap is a genuine dialetheia,
(b) it does not reduce to falsity, and (c) it is a fixed point of the
ENGAGR$\to$FSPLIT$\to$FFUSE cycle. The Frobenius condition $\mu \circ \delta =
\text{id}$ is satisfied on $B$ by construction.

\subsection{The unified type}

All eight problems encode to the following type, with one exception:

\[
\langle \text{Ð}_{\text{ω}};\ \text{Þ}_{\text{O}};\ \text{Ř}_{\text{=}};\ \text{Φ}_{\text{}};\ \text{ƒ}_{\text{ż}};\ \text{Ç}_{\text{@}};\ \text{Γ}_{\text{β}};\ \text{ɢ}_{\text{ˌ}};\ \text{⊙}_{\text{ÿ}};\ \text{Ħ}_{\text{!}};\ \text{Σ}_{\text{S}};\ \text{Ω}_{\text{z}} \rangle
\]

This is tier $\text{O}_{\text{inf}}$, crystal address $6\,738\,803$, cell~155
of the 400-cell crystal. The type is shared by:

\[
\begin{array}{lcc}
\textbf{Problem} & \textbf{Crystal Address} & \textbf{Chirality} \\ \hline
\text{Riemann Hypothesis}          & 6\,738\,803 & \text{Ħ}_{\text{!}} \\
\text{Yang-Mills mass gap}         & 6\,738\,803 & \text{Ħ}_{\text{!}} \\
\text{Hodge conjecture}            & 6\,738\,803 & \text{Ħ}_{\text{!}} \\
\text{Navier-Stokes regularity}    & 6\,738\,803 & \text{Ħ}_{\text{!}} \\
\text{P vs NP}                     & 6\,738\,803 & \text{Ħ}_{\text{!}} \\
\text{Birch--Swinnerton-Dyer}      & 6\,738\,800 & \text{Ħ}_{\text{A}} \\
\text{Odd Perfect Numbers}         & 6\,738\,803 & \text{Ħ}_{\text{!}}
\end{array}
\]

The convergence is striking: seven of the eight problems collapse to the exact
same tuple. The type is independent of the problem's specific
mathematical content---the RH's explicit formula and the Navier-Stokes helicity
winding produce the same 12-dimensional signature. This is not an artifact of
coarse granularity. The crystal has $17\,280\,000$ distinct types; $\text{O}_{\text{inf}}$
occupies only $1\,382\,400$ of them (8\%). That seven problems land in the same
cell (cell~155, one of 32 $\text{O}_{\text{inf}}$ cells) is a structural
prediction that the grammar makes and the data confirm.

\subsection{What the gap is, in each case}

The honest gap is different for each problem:

\begin{itemize}
\item {\bf RH}: The $\text{PM\_Z2}$ channel---Frobenius involution
      $\theta(s) = 1 - s$---fixes the critical line. The gap is unconditional
      proof that all non-trivial zeros satisfy $\text{Re}(s) = 1/2$.

\item {\bf YM}: Reflection positivity and area law imply a mass gap. The
      continuum limit $a \to 0$ is the gap: constructing the continuum quantum
      field theory from lattice regularisation.

\item {\bf Hodge}: The Grothendieck--Riemann--Roch factorisation
      $r = \text{cl} \circ \text{ch}$ is known. The Lefschetz $(1,1)$ theorem
      proves the Hodge conjecture for divisors. The gap is regulator surjectivity
      for $p \ge 2$.

\item {\bf NS}: Kato's theorem provides local-in-time smooth solutions.
      The Prodi--Serrin condition $(2/p + 3/q = 1)$ gives a regularity
      criterion. The gap is boundedness of the critical $H^{1/2}$ norm for all
      positive time.

\item {\bf P vs NP}: Three barriers---Baker--Gill--Solovay (relativisation),
      Razborov--Rudich (natural proofs), Aaronson--Wigderson (algebraisation)---
      each independently block a proof. The barrier triad joins to $B$.
      The gap is the $\text{O}_0 \neq \text{O}_1$ tier distinction.

\item {\bf BSD}: Modularity (Wiles~1995) and the functional equation
      $s \leftrightarrow 2 - s$ are established. Gross--Zagier and Kolyvagin
      prove rank $\le 1$. The gap is rank $\ge 2$.

\item {\bf OPN}: Euler's structure theorem ($N = p^\alpha m^2$,
      $p \equiv \alpha \equiv 1 \pmod 4$) constrains the form. The $2$-adic
      obstruction overdetermines the system. The gap is unconditional proof
      of nonexistence.
\end{itemize}

\section{The BSD Anomaly}

\subsection{What the crystal sees}

Six of the seven Clay problems---plus Odd Perfect Numbers---resolve to the
identical $\text{O}_{\text{inf}}$ tuple with chirality $\text{Ħ}_{\text{!}}$
(infinite Markov order). BSD sits three crystal addresses lower, at
$6\,738\,800$ instead of $6\,738\,803$, with chirality $\text{Ħ}_{\text{A}}$
(H2, persistent two-step memory).

The three-address gap is entirely in the chirality primitive:

\[
\begin{array}{c|c|c}
\text{Primitive} & \text{RH, YM, Hodge, NS, PvsNP, OPN} & \text{BSD} \\ \hline
\text{Ħ} & \text{Ħ}_{\text{!}}\ (\text{H}_\infty) & \text{Ħ}_{\text{A}}\ (\text{H2})
\end{array}
\]

In the Frobenius encoding, each notch of $\text{Ħ}$ contributes one unit to the
crystal address. The offset $\text{ord}(\text{Ħ}_{\text{!}}) - \text{ord}(\text{Ħ}_{\text{A}})
= 3 - 2 = 1$ moves through the positional weight of $\text{Ħ}$ in the address
computation. The result is exactly three addresses of difference.

This is not a discretisation artifact. It is a structural prediction that the
grammar extracts from the mathematical content of the problems.

\subsection{Why BSD is different}

The reason is visible in the mathematics itself:

\begin{itemize}
  \item {\bf Gross--Zagier (1986)}: If an elliptic curve $E/\mathbb{Q}$ has
        analytic rank $1$, then the Heegner point $P_K$ generates a subgroup of
        finite index in $E(\mathbb{Q})$, and the rank of $E(\mathbb{Q})$ is $1$.
  \item {\bf Kolyvagin (1989)}: If $L(E, s)$ does not vanish at $s = 1$, then
        the Selmer group is finite. Combined with Gross--Zagier: for analytic
        rank $\le 1$, $\text{rank}(E(\mathbb{Q})) = \text{ord}_{s=1}\, L(E, s)$.
\end{itemize}

These two results together constitute a proof of BSD for rank $\le 1$. The
proof is finite and complete. It does not rely on any unresolved conjecture.

This boundedness is captured by the chirality primitive. An $\text{H}_\infty$
system has unbounded memory depth---the recursion that generates its structure
does not terminate after a fixed number of steps. An H2 system has a maximum
memory depth of two: the structural recursion terminates after two alternating
steps. Gross-Zagier and Kolyvagin provide exactly this termination: the proof
establishes the result up to rank $1$, and the rank-$\ge 2$ case requires
genuinely new mathematics.

The other six problems have no such bounded proof component. Their partial
results---zero-density estimates for RH, reflection positivity for YM, Kato
local existence for NS, the three barriers for PvsNP---each approach the gap
but do not close it at any finite rank. The chirality is genuinely infinite.

\section{Discussion}

\subsection{What the work does and does not claim}

The dialetheic resolution presented here makes three precise claims:

First: each problem's honest gap can be formally encoded as a
Belnap-$\text{B}$ dialetheia in a kernel that does not explode on contradiction.
This is a formal theorem, proved by {\tt by decide} on every gap, in running
code.

Second: the type of each such encoding is tier $\text{O}_{\text{inf}}$
under the Imscribing Grammar. The convergence of seven problems to the same
tuple---and the BSD anomaly's three-address offset---are consequences of the
deterministic imscription procedure, not of subjective assignment.

Third: the kernel modification is sound and minimal. Three C++ files are
changed; the total new code is under 50 lines. The system remains fully
backward-compatible: when paraconsistent mode is off, the kernel behaves
identically to stock Lean 4.28.0.

The open question is what the crystal's structural classification implies
about the difficulty of the problems themselves. The $\text{O}_{\text{inf}}$
classification is a consequence of encoding the {\it gap as a dialetheia}. It
does not prove that the gap can be closed classically. The BSD anomaly is
informative precisely because BSD's gap {\it has} been partially closed by
Gross-Zagier and Kolyvagin, and the crystal detects this, whereas the other
gaps have not been similarly bounded.

\subsection{What remains to be done}

Several directions are natural:

\begin{itemize}
  \item {\bf The full paraconsistent proof theory already exists.} A prior
        version of this manuscript stated that a full paraconsistent proof theory
        was not attempted here. That was incorrect. The {\tt priests-engine}
        system---a complete ParaASM virtual machine with Belnap FOUR semantics,
        REPL, formal Lean verification, and an indefinitely-running Frobenius
        kernel---is exactly that theory. What was described as ``a deeper
        modification of the kernel's term language'' already exists as the Portal
        bridge between p4rakernel and {\tt priests-engine}: the
        ParaconsistentShell supports MEET/JOIN/TENSOR modes for bidirectional
        structural IPC, and {\tt priests-engine} proves $\mu \circ \delta =
        \text{id}$ on all four Belnap values.

        The natural next step is therefore not the construction of a new proof
        theory, but the integration of the two existing systems. A Lean tactic
        that invokes the ParaASM VM to verify dialetheic invariants, or a deep
        embedding of ParaASM as a Lean DSL, would unite the type-theoretic and
        operational sides of the same fact.

  \item {\bf Cross-crystal comparison.} The BSD anomaly shows that the crystal
        detects bounded proof depth. A systematic survey of problems in
        number theory, algebraic geometry, and analysis---mapping their
        chirality distributions---could reveal whether bounded chirality
        correlates with the existence of finite partial proofs.

  \item {\bf Dialogue with classical proof assistants.} The Lean 4 ecosystem
        is large and actively maintained. The paraconsistent fork is not a
        replacement but a complement. A natural next step is a tactic or plugin
        that allows switching between kernels within a single development.
\end{itemize}

\subsection{Closing the loop}

The introduction described the principle of explosion as a design decision, not
a logical necessity. The kernel modifications in Section~3 showed that this
decision can be unmade in under 50 lines of C++. The Belnap semantics in
Section~4 showed what replaces it: a four-valued logic in which contradictions
are contained. The Millennium resolutions in Section~5 showed that the
resulting framework classifies all eight problems at $\text{O}_{\text{inf}}$.
The BSD anomaly in Section~6 showed that the classification is not trivial---
the crystal discriminates between problems with bounded and unbounded chirality.

None of this proves the Riemann Hypothesis or resolves P vs NP. What it does
is relocate the boundary between what is known and what is open. The gap is no
longer a missing proof that must be supplied or else the system collapses. It
is a dialetheia that can be contained, named, and manipulated. The kernel
enforces the containment. The crystal records the type. The imscription is the
act of drawing the boundary.

And the paraconsistent proof theory that makes this possible is not a future
development. It is already running---25 billion paradoxes into its indefinite
loop, in a ParaASM machine that shares the same Belnap semantics, the same
Frobenius identity, and the same structural type as the kernel that contains it.
p4rakernel and {\tt priests-engine} are two sides of the same dialetheic coin:
one at the type-theoretic level, one at the operational level, both at
$\text{O}_{\text{inf}}$.

The principle of explosion was a decision. The decision not to explode is also
a decision. Both are implemented in running code. The difference is structural,
and the grammar measures it.

\begin{thebibliography}{10}
\bibitem{belnap1977}
N.~D.~Belnap, ``How a computer should think,'' in {\it Modern Uses of
  Multiple-Valued Logic}, 1977.

\bibitem{priest2006}
G.~Priest, {\it In Contradiction: A Study of the Transconsistent}, 2nd ed.,
Oxford University Press, 2006.

\bibitem{moura2015}
L.~de~Moura, S.~Ullrich, ``The Lean 4 Theorem Prover and Library,''
{\it CADE}, 2015.

\bibitem{grosszagier1986}
B.~Gross, D.~Zagier, ``Heegner points and derivatives of $L$-series,''
{\it Inventiones mathematicae}, 1986.

\bibitem{kolyvagin1989}
V.~Kolyvagin, ``Finiteness of $E(\mathbb{Q})$ and
$\text{SH}(E/\mathbb{Q})$ for a subclass of Weil curves,''
{\it Mathematics of the USSR-Izvestiya}, 1989.
\end{thebibliography}

\end{document}
